\noindent
This section summarises the extent to which the project objectives were achieved and identifies the most important directions for future development.

\subsection{Conclusions}

This project presented the design, implementation, and evaluation of a low-cost, modular haptic rendering system integrating a two-degree-of-freedom input device with a custom real-time simulation and visualisation framework.
The system successfully achieved its primary objective of enabling bidirectional interaction between a user-operated device and a virtual environment, with real-time computation of interaction forces and feedback through actuator torque control. The architecture was designed to prioritise modularity and extensibility, allowing clear separation between hardware, communication, simulation, and rendering subsystems.

Experimental evaluation showed that selected local subsystems could operate at nominal 1\,kHz rates, including embedded packet generation and host-side haptic/device-loop execution, with low host-side processing latency. However, the complete bidirectional haptic interaction should not be claimed as a true 1\,kHz end-to-end loop, because force computation operated on latest-available snapshots affected by communication delay and asynchronous simulation timing. The dominant limitation was the 15.60--15.67\,ms round-trip communication delay, which bounded achievable virtual stiffness.
Within this constraint, bounded and controllable contact interaction was observed during the tested interactions through a combination of architectural and control strategies, including snapshot-based communication and damping-based stabilisation. The resulting haptic feedback was usable and controllable, but not rigid or quantitatively force-calibrated, due to finite actuator torque, the absence of current sensing, belt slip under high load, and the limitation to single-axis actuation during final validation.

The project outcomes address the objectives defined in Section~\ref{sec:introduction}. A sub-£200 haptic interface was constructed with joint sensing and motor command capability, although final actuation was limited to one joint after actuator failure. The system demonstrated nominal 1\,kHz packet streaming while separately quantifying packet loss, out-of-order events, burst-like delivery, RTT, and host-side processing latency. Logged simulated rendered-force data gave an effective stiffness of 489.57\,N\,m$^{-1}$, within 5\% of the configured 500\,N\,m$^{-1}$ value, with $R^2=0.9973$. Saturation, torque limiting, slew-rate limiting, watchdog behaviour, motor-enabled contact, and constrained-motion tests were also validated, while documenting the communication, actuation, sensing, and mechanical limitations that affect generalisation.

Overall, the project demonstrates that bounded and usable haptic interaction can be achieved using low-cost hardware and a modular software architecture, even in the presence of significant communication latency. The main insight is that high-rate local control and an accurate command-side force model are not sufficient on their own: communication delay, actuator capability, transmission design, and the absence of a persistent high-rate local contact-state model jointly determine achievable haptic fidelity.

\subsection{Future Work}
A key priority is reducing communication latency. The current ST-Link virtual COM pathway appears to introduce buffering and scheduling delays that dominate the closed-loop response, so future work should compare lower-latency alternatives such as direct UART via a dedicated USB--UART bridge, native USB bulk transfer, CAN bus, or Ethernet-based communication.

Future work should extend the existing multi-rate architecture so that the haptic loop maintains a reduced local model of currently contacted objects at 1\,kHz, while slower physics, rendering, or application layers remain responsible for broader scene evolution and consistency. The high-rate loop would continuously evolve the proxy, contact manifold, and local object interaction state, reducing dependence on low-rate world snapshots and improving contact continuity, following the motivation of multilevel haptic rendering for complex environments \cite{ruspiniHapticDisplayComplex1997}.

A related extension would be to move the haptic rendering loop, or at least the local contact/proxy evolution, onto the embedded controller. Communication with the host would then focus on object interaction parameters, local contact primitives, or periodic scene-state updates, rather than every force update depending on a delayed host--device exchange. This would lower effective actuator-side latency, with the trade-off that the embedded model would need to be simplified using local planes, SDF patches, virtual fixtures, or reduced contact geometry.

On the control side, incorporating task-space dynamics using an operational-space formulation \cite{khatibUnifiedApproachMotion1987a} would allow the manipulator's inertia and dynamics to be compensated during interaction, improving force-rendering transparency.

Hardware improvements are also important. Current sensing, together with external force measurement, would enable more accurate torque control, quantitative physical force validation, and more advanced impedance or operational-space control. Transmission design could be improved by replacing the slipping GT2 belt stage with a lower-backlash alternative such as a capstan drive, while higher-torque actuators and restored two-degree-of-freedom actuation would improve interaction realism.

Further improvements to the haptic rendering model could address the force-direction inconsistencies observed during large penetration. Constraint-based contact solvers could enforce continuous proxy motion across timesteps, reducing surface-normal discontinuities caused by re-projecting the proxy from scratch. Contact entry and exit transitions should also be smoothed by blending proxy velocity state across release and discarding stale torque commands before they affect the user.

From a software and evaluation perspective, integration with platforms such as NVIDIA Isaac Sim or MATLAB/Simulink would enable more realistic robotic scenarios, while quantitative force measurement, user studies, and task-based metrics would better assess effectiveness as a training tool.

Together, these developments would enable the system to evolve from a low-cost prototype into a more capable and general-purpose haptic interface platform for robotic interaction and training applications.
